\section{评估与性能结论：比较前必须统一分母}\label{sec:evaluation}
\paper{第4.1节、图3--4（第6--7页），附录F.1--F.2（第17--18页）}
\subsection{评估脚本具体做了什么}
\code{evaluate.py} 每个 clip 分配一个环境，全部从该 clip 首帧参考状态启动。评估关闭初始扰动、关节零位与质心随机化，摩擦固定1.0，不加观测噪声和动作噪声。使用 checkpoint 的策略及其归一化统计，推理目标仍是10帧未来参考；\code{pi.eval()} 不等于 \code{no_grad()}，脚本另用装饰器显式关闭梯度。

环境达到硬失败或 clip 结束后，将该环境标为完成；脚本为批处理仍让其物理前进，但累计回报和失败统计只在未完成 mask 下更新，并限制完成环境的帧索引避免越界。这不是失败后 reset 再继续补足同一个 clip 的成绩。

设第 $i$ 个 clip 重采样后有 $L_i$ 帧，回报只累计首次失败/结束前的 $T_i$ 个转移。归一化为
\begin{equation}
 \bar G_i=\frac{\sum_{t=0}^{T_i-1}r_{i,t+1}}{(L_i-1)\,r_{\max}},\qquad r_{\max}=0.1.
\end{equation}
分母始终是\textbf{原 clip 完整转移数}，而不是活过的 $T_i$。这样提前失败后没有继续获得的奖励也反映在分数里。该回报无时间折扣，且包含负正则项，因此不是纯姿态误差或成功率，也不保证下界为0。
\source{\src{supermimic/evaluate.py}{52}{96}；\src{supermimic/env.py}{64}{66}；\src{supermimic/env.py}{215}{223}}

脚本同时输出两种聚合：
\begin{align}
 \texttt{norm\_return}&=\frac1C\sum_{i=1}^{C}\bar G_i,\\
 \texttt{norm\_return\_framew}&=
 \frac{\sum_i\sum_t r_{i,t+1}}{\sum_i(L_i-1)r_{\max}}.
\end{align}
前者每个 clip 等权，后者每个参考转移等权，长 clip 权重更高。论文附录F.1明确了首帧、摩擦、无噪声与失败判据，但没有在该处完全展开这两种聚合的选择。仓库自报数字应按它实际写入的 \code{norm_return} 理解；不能在不同汇总分母下直接宣布严格复现同一曲线。

\subsection{论文原图：规模扩展与 LAFAN1 对照}
\begin{figure}[H]
\centering
\includegraphics[width=\textwidth]{figures/paper_fig3_scaling.pdf}
\caption{论文原图3：跟踪回报、训练成本与数据规模。引用自 Liao 等（2026）图3，第7页\cite{liao_supermimic_2026}，从 arXiv:2610.28480v1 的 PDF 矢量裁切。原图(a)比较 SONIC 与 SuperMimic 的 GPU-hours；(b,c)展示1至264小时动作库的成本与8步预测误差。保留原坐标轴、图例与预算标记；不是40段 LAFAN1 仓库的训练曲线。}
\label{fig:paper-scaling}
\end{figure}

图\ref{fig:paper-scaling}(a)的主要基准是129,514段 BONES-SEED 动作；(b,c)刻画同一机器人动力学在不同数据覆盖下的扩展行为。它支持“一个模型可共享于很多动作”的论文动机，但当前代码没有提供这一数据扩展实验的完整复现链路。特别是(c)的 nominal-error 归一化与代码 \code{wm_err8} 的自适应标准化不能直接视为同一标尺，硬件与训练数据规模的差异也应一并保留。

\begin{figure}[H]
\centering
\includegraphics[width=\textwidth]{figures/paper_fig4_controlled.pdf}
\caption{论文原图4：LAFAN1 受控比较与 critic/horizon 消融。引用自 Liao 等（2026）图4，第7页\cite{liao_supermimic_2026}，从 arXiv:2610.28480v1 的 PDF 矢量裁切。所有子图使用40段、约2.45小时的 LAFAN1 动作；(a)是 checkpoint 在完整动作库上的评估回报，(b,c)使用训练回报。阴影标出默认 $H=32$，叉号表示未达目标。}
\label{fig:paper-controlled}
\end{figure}

图\ref{fig:paper-controlled}(a)才是该仓库 README 声明的直接对照目标；(b,c)则解释了默认 \code{--horizon 32} 及末端 critic 的设计背景。代码提供 \code{--no_critic} 和可配置 horizon，但这不等于已交付完整消融结果。还需注意：训练 episode 采用基于回报的重置分布，而(a)从每个 clip 首帧评估，所以\textbf{不能把(b,c)的训练回报数值与(a)的评估分数逐点等同}。图中成绩均归原论文，仓库作者记录另在下表单独列出。

\subsection{三个层次的成绩不应混在一张结论里}
\begin{table}[H]
\centering\small
\noindent\begin{tabularx}{\textwidth}{@{}P{30mm}P{40mm}Y@{}}
\toprule
来源/设置 & 报告的训练成本 & 指标与适用范围\\
\midrule
论文LAFAN1，图4(a) & 1张H100，1.2h & 40段、约2.45h；归一化回报达到0.82\\
仓库LAFAN1，里程碑 & 1张RTX PRO 6000，1.24h & 迭代12,000；回报0.824；仍有1/40失败\\
仓库LAFAN1，里程碑 & 同GPU，1.47h & 迭代14,000；回报0.829；0/40失败\\
论文BONES-SEED & 1张H100，3h & 129,514段、264.56h；25段失败\\
论文BONES-SEED & 2张H100，9.8h & 同主基准；0段失败；成本19.6 GPU-h\\
本次报告 & 未进行训练 & 16项CPU检查通过；不产出策略回报或失败数\\
\bottomrule
\end{tabularx}
\caption{论文成绩与非官方仓库自报成绩分开列示。仓库数字来自 README 与 MILESTONES，不是本次测量。}
\end{table}
\source{\src{README.md}{1}{24}；\src{MILESTONES.md}{1}{23}}

论文将 SuperMimic 的3 GPU-h与 SONIC 的 $128\times168=21,504$ GPU-h比较，数量级比约为7,168；但论文也说明GPU型号、训练数据规模和可用峰值算力不同。该计算是理解论文预算口径，不是这个40段仓库已重新测得“千倍加速”。同理，作者1.47小时是单GPU训练时间记录，不包含全部构建、下载和评估成本，也不能当作任意硬件的必达承诺。

\subsection{成功率、跟踪质量与泛化是三种问题}
失败判据只有高度和重力z分量差。\textbf{0/40失败不表示每一关节都准确匹配}，也不表示没有任何物理异常；还需看回报、分项误差、动作平滑和具体动作轨迹。

论文附录F明确，其主基准是训练动作上的覆盖率，而非未见动作泛化。该仓库也在训练所用LAFAN1动作库上评估，不能把40/40成功写成“对未见动作100\%泛化”。速度接口任务和实机迁移属于另一层证据，当前仓库没有提供对应完整复现实验。

\subsection{当前公开材料足以支持哪些判断}
固定提交包括结果曲线图、演示GIF和文字里程碑，但没有随版本跟踪完整训练CSV与checkpoint。因此可以审查实现是否具有论文所述机制，亦可准确转述作者的结果记录；不能只凭图像和README独立审计所有训练时长、每段动作指标或训练随机性。

若进一步做严格实验，应固定数据文件校验和与资产转换版本，记录实际命令、GPU、驱动、CUDA 与 PyTorch、碰撞开关和精度，保留每clip回报/失败信息，并分别报告clip-mean与frame-weighted分数。若要声称观测严格对齐，应先单独消融新增6维参考朝向，而不是悄悄删除它后仍沿用已有成绩。
